\section*{Chapter 1: Introduction}
\setcounter{section}{1} \setcounter{subsection}{0}
\phantomsection \addcontentsline{toc}{section}{\textit{Chapter 1: Introduction}}

This report presents the design, implementation and evaluation of an LLM multi-agent negotiation framework for 5G network slice resource management. The framework combines a two-tier scheduling architecture with specialised LLM agents that collaborate through a structured proposal--evaluation--negotiation protocol to allocate shared radio bandwidth among competing slice types. This chapter introduces the research context, states the problem addressed, defines the research objectives and scope, summarises the contributions, and outlines the organisation of the remainder of this report.

\subsection{Background and Motivation}

Fifth-generation (5G) network slicing enables the creation of multiple logical networks over a shared physical infrastructure, with each slice tailored to a distinct class of service~\cite{afolabi2018slicing}. The three canonical slice types defined by the 3rd Generation Partnership Project (3GPP) --- Enhanced Mobile Broadband (eMBB), Ultra-Reliable Low-Latency Communication (URLLC) and massive Machine-Type Communication (mMTC) --- differ substantially in their throughput, latency, reliability and connection-density requirements. When these slices share a common pool of radio resources, their Quality of Service (QoS) targets become inherently competitive: heavy bandwidth consumption by eMBB users streaming high-definition video may threaten the latency guarantees required by URLLC users, while periodic mass reporting by mMTC devices can reduce the access resources available to other slices~\cite{afolabi2018slicing}.

Traditional approaches to network slice resource management include rule-based heuristic allocation and Deep Reinforcement Learning (DRL)-based optimisation~\cite{li2018drl,sciancalepore2019rlnsb}. While these methods can encode priorities and constraints explicitly, they generally struggle to integrate Service Level Agreement (SLA) constraints, operator-defined business priority policies and dynamically evolving user behaviour patterns into a unified and adaptive decision-making framework. Rule-based methods generalise poorly to traffic conditions that deviate from those anticipated at design time, whereas DRL methods require extensive training data and careful reward engineering, and provide limited interpretability of their allocation decisions.

The emergent reasoning capabilities of Large Language Models (LLMs) present a potential alternative~\cite{zhou2024llmtelecom}. Since 2023, researchers have begun exploring LLMs as intelligent decision engines for network management. Frameworks such as WirelessAgent~\cite{tong2025wirelessagent} have demonstrated that a cognitive architecture combining perception, planning and action, supported by tool calling, can approach rule-optimal performance in slicing scenarios without task-specific training. However, existing LLM-based approaches predominantly adopt a centralised single-agent design in which one model is responsible for resource decisions across all slices, without exposing explicit negotiation among slice interests. This observation motivates the present work.

\subsection{Problem Statement}

The core problem addressed by this project is the dynamic allocation of shared radio bandwidth among multiple 5G network slices under heterogeneous SLA constraints. Two specific limitations of existing approaches are identified:

\begin{enumerate}
    \item Existing LLM-based network management frameworks either adopt a centralised single-agent design that does not support explicit inter-slice negotiation~\cite{tong2025wirelessagent,shao2024wirelessllm}, or they propose multi-agent architectures whose validation is confined to tasks other than slice resource allocation~\cite{jiang2024commllm}. A complete, end-to-end multi-agent LLM framework dedicated to network slice resource management and validated in a reproducible simulation environment has not yet been demonstrated.

    \item General-purpose multi-agent LLM frameworks such as MetaGPT~\cite{hong2023metagpt} and ChatDev~\cite{qian2023chatdev} provide rich coordination mechanisms, but these are designed for predominantly cooperative tasks. Network slice resource allocation is better modelled as a constrained multi-party optimisation problem in which the objectives of different slices compete under a fixed capacity budget~\cite{caballero2019games}. Structured negotiation and conflict-resolution protocols tailored to this competitive setting remain underexplored.
\end{enumerate}

\subsection{Research Objectives}

This project aims to design and implement an LLM multi-agent negotiation-based resource allocation framework in a 5G network slicing simulation environment, evaluate it against established baselines, and assess the contribution of the negotiation mechanism. The research is guided by two questions:

\textbf{RQ1:} Can a multi-agent LLM framework match or exceed the performance of traditional rule-based strategies and DRL methods in network slice resource management, as measured by SLA satisfaction rate, resource utilisation and dynamic load adaptability?

\textbf{RQ2:} Does the structured negotiation mechanism between agents provide measurable advantages over a centralised single-round allocation by the coordinator alone, particularly in terms of inter-slice fairness and SLA satisfaction under load conflict scenarios?

\subsection{Scope and Assumptions}

The following boundaries define the scope of this project:

\begin{itemize}
    \item The evaluation is conducted entirely in simulation using a custom Gymnasium-based environment. No real network hardware or live traffic traces are used.
    \item The system considers three slice types (eMBB, URLLC, mMTC) sharing a total bandwidth of 100\,MHz, with a maximum of 100 concurrent users.
    \item SLA definitions are simulation-layer approximations that capture relative priority isolation rather than strict 3GPP end-to-end physical-layer metrics. Specifically, eMBB requires average user throughput $\geq$ 50\,Mbps, URLLC requires 99th-percentile queuing delay $\leq$ 10\% of the eMBB delay, and mMTC requires access success rate $\geq$ 95\%.
    \item The LLM backbone used is a commercial API (GPT-4 series) with temperature fixed at~0 to minimise stochastic variability.
    \item Each experimental condition is evaluated over 5 independent random seeds across two traffic scenarios (steady-state and burst).
\end{itemize}

\subsection{Contributions}

The main contributions of this project are:

\begin{enumerate}
    \item \textbf{A reproducible 5G network slicing simulation environment} built on the Gymnasium framework~\cite{brockman2016gym}, with clearly defined state space (15 dimensions), action space (3-dimensional continuous allocation), reward function and SLA evaluation logic. All parameters and interfaces are publicly available.

    \item \textbf{A two-tier scheduling LLM multi-agent decision framework} that reconciles the inference latency of LLMs with the real-time demands of radio resource management. The upper tier features multiple specialised LLM agents --- three slice managers and one global coordinator --- that collaborate through a structured proposal--evaluation--negotiation protocol. The lower tier executes cached rule-based strategies at millisecond granularity. Two negotiation strategies, Priority Arbitration and Proportional Compromise, are implemented and compared.

    \item \textbf{A systematic empirical evaluation} comparing six methods (two rule-based baselines, one DRL baseline, one single-agent LLM baseline and two multi-agent LLM variants) together with a negotiation ablation study, across two traffic scenarios and five random seeds per condition.
\end{enumerate}

\subsection{Report Organisation}

The remainder of this report is structured as follows. Chapter~2 reviews the technical background and related work, covering network slicing architectures, SLA-aware scheduling, LLM-based agent frameworks for network management, multi-agent coordination mechanisms, and the identified research gaps. Chapter~3 presents the design and implementation of the proposed framework, including the system architecture, two-tier scheduling, agent definitions, negotiation protocol, simulation environment, LangGraph-orchestrated workflow and baseline methods. Chapter~4 describes the experimental setup, presents the results of the two main experiments, and discusses findings, limitations and practical implications. Chapter~5 concludes the report with a summary of achievements, a personal reflection, and suggestions for further work.

%% ======== TEMPLATE GUIDANCE (retained per instructions) ========
% The \textit{Introduction} is one of the most important parts of your report as it gives a brief overview that will make the reader understand (i) what you set out to do and (ii) what you achieved. Many readers will read the \textit{Introduction} first and then the  \textit{Conclusions} to get an overview before reading the detail – so it is important that both sections are very carefully written.
%
% It is very important that \textit{Introduction} introduces the \textbf{report} and the \textbf{project} – it is not there to introduce the subject in general.
%
% There are no rules on how to write an introduction, but it should include what the project is about, give \textit{a very short description} of the technical context in which the project is carried out and explain the motivation for the work. If you are doing an implementation project it must explain what functionality the system \textbf{realises}, and if you are doing a research project what is the novelty of the approach used.
%
% Very importantly, it should clearly indicate what you have done for the project.
%
% A good introduction should be no more than 4 pages.
